\chapter{Manual de instalare și utilizare}
\label{ch:instalare}
\pagestyle{fancy}

Acest capitol descrie modul în care sistemul trebuie asamblat, conectat și pus în
funcțiune, precum și modurile în care poate fi utilizat. Sunt prezentate pe rând
componentele necesare, alimentarea, conectarea plăcilor între ele și, în final,
comenzile prin care utilizatorul controlează mâna robotică.

\section{Componente și resurse necesare}
\label{sec:manual-componente}

Pentru punerea în funcțiune a sistemului sunt necesare placa FPGA Basys3, mâna robotică
asamblată împreună cu cele opt servomotoare, cele două izolatoare ADuM1400 pentru
semnalele de comandă, izolatorul SPI MAX14850, convertorul analog-digital AD7124-8,
regulatorul liniar ADP165 și modulul de alimentare DC2468A. Pe lângă acestea, este
nevoie de o sursă de alimentare externă și, opțional, de un calculator gazdă pentru
modul de urmărire în timp real.

\section{Alimentarea sistemului}
\label{sec:manual-alimentare}

Alimentarea întregului sistem pornește de la o sursă externă conectată la modulul
DC2468A. Aceasta trebuie să furnizeze o tensiune de 12~V și un curent de cel puțin 8~A,
deoarece cele opt servomotoare, atunci când acționează simultan sub sarcină, pot cere
un curent considerabil. În locul unei surse de laborator, se poate folosi și un
încărcător suficient de puternic, adaptat corespunzător pentru a furniza tensiunea și
curentul necesare.

Modulul DC2468A coboară tensiunea de la 12~V la 5~V, tensiune folosită pentru
alimentarea servomotoarelor. Tot din DC2468A este alimentat și regulatorul liniar
ADP165, care coboară mai departe tensiunea la 3,3~V pentru convertorul analog-digital.
De asemenea, latura izolată a izolatoarelor, cea dinspre servomotoare și convertor,
este alimentată tot din DC2468A, astfel încât întreaga parte de putere folosește aceeași
sursă.

Partea logică, în schimb, este alimentată separat, direct de la placa Basys3. Aceasta
furnizează tensiunea de 3,3~V pentru latura dinspre placă a izolatoarelor, adică pentru
partea care primește semnalele digitale de la FPGA. În felul acesta, cele două domenii
de alimentare, cel logic dinspre placă și cel de putere dinspre sursă, rămân complet
separate galvanic, alimentarea provenind din surse diferite pe cele două laturi ale
izolatoarelor.

\section{Conectarea plăcilor}
\label{sec:manual-conectare}

Placa Basys3 trimite cele opt semnale PWM prin conectoarele sale Pmod. Primele patru
canale, corespunzătoare primelor patru unghiuri din cadru, sunt scoase prin conectorul
JA, iar următoarele patru, corespunzătoare celorlalte unghiuri din cadru, prin
conectorul JB. Fiecare dintre aceste semnale intră într-unul dintre cele două
izolatoare ADuM1400, care realizează atât izolarea galvanică, cât și ridicarea nivelului
logic de la 3,3~V la 5~V. Corespondența dintre canale, pinii fizici ai plăcii și
articulațiile comandate este dată în Tabelul~\ref{tab:pini-pwm}.

\begin{table}[H]
\centering
\caption{Corespondența pinilor PWM de pe placa Basys3.}
\label{tab:pini-pwm}
\begin{tabular}{c l l l}
\toprule
\textbf{Canal} & \textbf{Pmod} & \textbf{Pin FPGA} & \textbf{Articulație} \\
\midrule
1 & JA1 & J1 & Deget mic \\
2 & JA2 & L2 & Inelar \\
3 & JA3 & J2 & Mijlociu \\
4 & JA4 & G2 & Arătător \\
5 & JB1 & A14 & Police (palmar) \\
6 & JB2 & A16 & Police (tensiune) \\
7 & JB3 & B15 & Încheietură \\
8 & JB4 & B16 & Cot \\
\bottomrule
\end{tabular}
\end{table}

Se folosesc numai primii patru pini de semnal ai fiecărui conector, restul pinilor
Pmod rămânând neconectați. Pe lângă pinii de semnal, fiecare conector Pmod pune la
dispoziție și pini de alimentare de 3,3~V și de masă, folosiți pentru alimentarea
laturii dinspre placă a izolatoarelor.

Primul izolator ADuM1400 preia primele patru canale din cadru, iar al doilea preia
următoarele patru. Pe fiecare izolator, canalele sunt organizate identic pe cele două
laturi, în ordinea alimentare, masă și cele patru canale de semnal, notate A, B, C și D.
Astfel, atât intrarea dinspre placă, cât și ieșirea dinspre servomotoare urmează aceeași
dispunere a pinilor, ceea ce simplifică conectarea. Latura dinspre placă a fiecărui
izolator este alimentată de la Basys3, iar latura dinspre servomotoare de la sursa de
putere, prin DC2468A, păstrând astfel separarea galvanică. Pentru ca izolarea să fie
efectivă, este important ca alimentările și masele celor două laturi să fie conectate
corect și să nu fie amestecate între ele.

Pe calea de feedback, convertorul AD7124-8 comunică cu placa prin izolatorul SPI
MAX14850. Semnalele SPI sunt scoase de pe placă prin conectorul JC, conform
Tabelului~\ref{tab:pini-spi}. Regulatorul ADP165 se conectează la pinii de alimentare
ai izolatorului MAX14850 de pe latura convertorului, pentru a alimenta atât izolatorul,
cât și convertorul cu tensiunea de 3,3~V.

\begin{table}[H]
\centering
\caption{Corespondența pinilor SPI de pe conectorul JC al plăcii Basys3.}
\label{tab:pini-spi}
\begin{tabular}{c l l}
\toprule
\textbf{Semnal SPI} & \textbf{Pmod} & \textbf{Pin FPGA} \\
\midrule
CS (selecție) & JC1 & K17 \\
MOSI (date către ADC) & JC2 & M18 \\
MISO (date de la ADC) & JC3 & N17 \\
SCLK (ceas) & JC4 & P18 \\
\bottomrule
\end{tabular}
\end{table}

Semnalele SPI trec mai întâi prin izolatorul MAX14850, care le transferă peste bariera
de izolare către convertor, legarea făcându-se după numele semnalelor, deoarece ordinea
pinilor diferă între cele două plăci. Ansamblul complet al conexiunilor dintre placă,
izolatoare, convertor și servomotoare este prezentat în Figura~\ref{fig:manual-conexiuni}.

\begin{figure}[t]
	\centering
	\includegraphics[width=\textwidth]{figs/schema_conexiuni.pdf}
	\caption{Schema conexiunilor dintre placa Basys3, izolatoare, convertor și
	servomotoare, cu alimentările aferente.}
	\label{fig:manual-conexiuni}
\end{figure}

\section{Moduri de utilizare}
\label{sec:manual-utilizare}

După conectarea și alimentarea corectă a sistemului, mâna robotică poate fi controlată
în trei moduri, selectate prin comutatoarele de pe placa Basys3. Modul activ este
determinat de starea acestor comutatoare, iar trecerea de la un mod la altul se face
simplu, prin schimbarea poziției lor.

În mod implicit, când comutatoarele dedicate sunt coborâte, sistemul funcționează în
modul normal, în care unghiurile sunt primite prin UART de la calculatorul gazdă. Acesta
este modul folosit pentru urmărirea în timp real a mâinii operatorului.

Comutatorul SW0 activează modul demonstrativ. În acest mod, mâna execută automat o
secvență predefinită de mișcări, fără a mai fi nevoie de conexiunea cu calculatorul,
fiind util pentru demonstrarea funcționării.

Comutatorul SW15 activează modul manual, în care fiecare articulație este comandată
individual de câte un comutator dedicat. Comutatoarele SW1 până la SW8 corespund celor
opt servomotoare, în ordinea din cadru, adică degetul mic, inelarul, mijlociul,
arătătorul, mișcarea palmară a degetului mare, tensiunea degetului mare, încheietura și
cotul. Un comutator ridicat trimite articulația corespunzătoare către poziția maximă,
iar unul coborât către poziția minimă, permițând astfel controlul direct al fiecărei
mișcări. Acest mod este util atât pentru verificarea individuală a servomotoarelor, cât
și pentru reglajele mecanice.

Butonul central al plăcii are rolul de resetare a sistemului, readucându-l în starea
inițială. Indiferent de modul ales, comenzile trec prin filtrul Kalman și prin
generatoarele PWM, astfel încât mișcarea rămâne lină în toate situațiile.
